\chapter{Systèmes auxiliaires}


Dans ce module, nous allons décrire les principaux équipements auxiliaires que l’on rencontre sur un site d’incinération des déchets ménagers. Par exemple, les équipements de production d’air comprimé assurent le fonctionnement des vannes de régulation ainsi que l’outillage pneumatique. 
\vspace{0.3cm}


Les équipements du réseau de vapeur assurent le transport, la purge, le sectionnement et la régulation de la circulation de la vapeur. D’autres équipements, les groupes électrogènes et les alimentations sans interruption, comme les onduleurs, assurent la production d’énergie électrique des principaux organes de sécurité, même en cas de panne du réseau EDF. D’autres équipements réservés à la lutte contre l’incendie ont également un rôle très important.

\vspace{0.3cm}

À la fin de ce module, vous pourrez identifier ces équipements, expliquer, leur utilité et leur principe de fonctionnement. Vous pourrez ainsi assurer au mieux l’entretien et la surveillance de base.


\section{Air Comprimé}


\subsection{Rôle de l'Air Comprimé}

L'air est utilisé comme fluide de travail pour les applications suivantes :
\begin{itemize}
\item  nettoyage, ramonage ;
\item  alimentation des servomoteurs des vannes de régulation ;
\item  alimentation d’outils, tels que burins, clés à chocs, etc. ;
\item  automatisation et robotisation de tâches répétitives.

\end{itemize}


L’air comprimé est disponible en grande quantité car il est pris à l’atmosphère. Il peut être stocké facilement. Il n’a pas besoin d’être recyclé, une fois utilisé il est rejeté à l’air libre.
\vspace{0.3cm}



Il convient parfaitement pour les travaux de nettoyage à l’intérieur des fours d’incinération.
Il peut être utilisé quelle que soit l’ambiance de travail, par exemple dans des atmosphères explosives. 




\subsection{Production d'air comprimé}

Un réseau de production d'air comprimé est composé des éléments suivants :

\begin{figure}[H]
\centering
\includegraphics[width=\textwidth]{Images/M21/M21_1.png}
\end{figure}

L’air comprimé est aspiré à l’extérieur par un réseau de conduit. Il est préfiltré pour enlever les poussières les plus grosses. Il est ensuite comprimé dans le compresseur. Il est stocké dans un réservoir. Il est ensuite filtré, déshydraté par le sécheur, puis à nouveau filtré avant d’être distribué jusqu’au lieu d’utilisation.

\subsection{Compresseurs}

\textbf{Rôle des compresseurs}
\vspace{0.3cm}

Le compresseur d’air doit assurer la production d’air sous un débit et une pression suffisante pour alimenter les équipements pneumatiques.
Il existe plusieurs manières de produire de l’air comprimé, par le biais d’unités décentralisées situées à proximité de l’utilisateur ou par le biais d’unités centralisées dans un local technique, qui nécessitent la pose de canalisations de distribution jusqu’à l’utilisateur.
\vspace{0.3cm}

Le compresseur présenté ci-dessus est une unité décentralisée. Il est transporté sur le lieu d’utilisation de l’air comprimé. Il est composé d’un compresseur à piston et d’un réservoir intégré. Il existe aussi des modèles pour lesquels le réservoir est séparé.
\vspace{0.3cm}

Le compresseur présenté ci dessus est une unité centralisée. Il est placé dans un local technique. L’air comprimé produit est acheminé par un réseau de canalisations.
\vspace{0.3cm}

Le débit des compresseurs est déterminé en fonction de la consommation d’air comprimé des outillages et de l’instrumentation.

\begin{figure}[H]
\centering
\includegraphics[width=0.8\textwidth]{Images/M21/M21_2.png}

\end{figure}

En règle générale, la pression d’utilisation des outils est de l’ordre de 7 bars. Le compresseur fournira donc l’air comprimé à une pression légèrement supérieure.
Il existe plusieurs types de compresseurs :

\begin{itemize}

\item compresseurs à pistons ou à vis,
\item compresseurs centrifuges. 
\end{itemize}

\textbf{Les compresseurs volumétriques}
\vspace{0.3cm}

Les compresseurs sont dits volumétriques car ils compriment une quantité d’air fixe à chaque tour de moteur. La compression est liée à la variation de volume d’une enceinte où est emprisonné l’air. 

\vspace{0.3cm}
\textbf{Compresseurs à pistons}
\vspace{0.3cm}

L’air est comprimé dans un cylindre par un piston. Ce dernier est animé d’un mouvement de va et vient grâce à un système de bielle et de vilebrequin. Le cylindre est équipé de clapets d’admission et de refoulement d’air.

\begin{figure}[H]
\centering
\includegraphics[width=0.4\textwidth]{Images/M21/M21_3.png}
\end{figure}


Dans le schéma ci-contre, on retrouve la phase d’aspiration et la phase de refoulement. 
\vspace{0.3cm}

À gauche, le piston descend, le clapet d’admission s’ouvre. L’air entre dans la chambre de compression. 
\vspace{0.3cm}

À droite, le piston remonte, le clapet d’admission se ferme tandis que le clapet de refoulement s’ouvre, l’air est refoulé vers l’extérieur.


\begin{figure}[H]
\centering
\includegraphics[width=0.6\textwidth]{Images/M21/M21_4.png}
\end{figure}



On rencontre de nombreux types de compresseurs à piston : le compresseur simple effet et le compresseur double effet.
Dans le compresseur simple effet, la compression n’a lieu qu’à une extrémité du cylindre. 
Dans le compresseur double effet, la compression a lieu aux deux extrémités du cylindre.
\vspace{0.3cm}

Lorsque la pression d’utilisation est de l’ordre de 7 bars, on utilise des compresseurs à deux étages, ou bi-étagés. L’air est comprimé dans le premier étage à une pression de 1 bar, puis il est envoyé dans un second étage où il est comprimé à sa pression finale.






\begin{figure}[H]
\centering
\includegraphics[width=0.5\textwidth]{Images/M21/M21_5.png}
\end{figure}




\textbf{Compresseurs à vis}
\vspace{0.3cm}

Le compresseur a deux vis s'entrecroisant. Lorsque les vis tournent, l'air est aspiré, comprimé et envoyé.


\begin{figure}[H]
\centering
\includegraphics[width=0.5\textwidth]{Images/M21/M21_6.png}
\end{figure}







\textbf{Compresseurs centrifuges}
\vspace{0.3cm}

Ils fonctionnent selon le même principe que les pompes centrifuges. Une roue à aubes tourne à grande vitesse. L’air, admis parallèlement à l’axe de rotation, en contact avec la roue, est entraîné en rotation à grande vitesse. Il est ensuite envoyé vers la canalisation de refoulement où l’énergie cinétique (due à la vitesse) est transformée en énergie de pression.

\begin{figure}[H]
\centering
\includegraphics[width=0.3\textwidth]{Images/M21/M21_7.png}
\end{figure}

\vspace{0.3cm}
\textbf{Fonctionnement des compresseurs}
\vspace{0.3cm }

Lorsque l’on comprime de l’air, il résulte un échauffement important qui peut être dommageable pour le compresseur si on ne le refroidit pas.

\vspace{0.3cm}

Le compresseur est refroidi par l’air pris dans le local où il est placé. Un système de batterie d’échange air/air permet de refroidir l’air comprimé ainsi que l’huile de lubrification des équipements du compresseur.

\vspace{0.3cm}

On surveillera donc les températures de l’air comprimé et de l’huile selon les recommandations du fabricant. Les températures sont visibles sur des indicateurs placés en façade du compresseur.

\begin{figure}[H]
\centering
\includegraphics[width=0.9\textwidth]{Images/M21/M21_8.png}
\end{figure}

\subsection{Les filtres}

\textbf{Rôle des filtres}
\vspace{0.3cm }

L'air atmosphérique est composé des éléments suivants :
\begin{itemize}

\item air sec
\item vapeur d'eau
\item poussière
\item polluants.
\end{itemize}

Sans aucun traitement de l’air, l’eau en se condensant corrode les canalisations en acier. Il se forme de la rouille, dommageable pour les équipements pneumatiques.
La présence des particules solides encrasse les équipements, diminue leurs performances et provoque une usure prématurée.
\vspace{0.3cm}


Il faut alors filtrer l’air pour éliminer les particules solides (rouille, calamine…) et l’eau condensée. Cette fonction est assurée par des filtres à air.
En bas de ces filtres, on trouve des systèmes de purges automatiques ou manuelles qui permettent de chasser les impuretés et les particules filtrées.


\begin{figure}[H]
\centering
\includegraphics[width=0.3\textwidth]{Images/M21/M21_9.png}
\end{figure}

\vspace{0.3cm}
\textbf{Filtre séparateur d'eau}
\vspace{0.3 cm}

L’air pénètre dans le filtre et prend un mouvement giratoire (de rotation). Les particules solides et liquides plus lourdes que l’air sont projetées contre les parois de la cuve et se déposent au fond. L’air passe ensuite au travers un matériau filtrant qui retient les impuretés.

\vspace{0.3cm}

L’eau condensée et les impuretés peuvent être purgées en bas de la cuve par un purgeur manuel ou par un purgeur automatique à flotteur.


\vspace{0.3cm}
\textbf{Filtres à huile}
\vspace{0.3cm}


L'huile ou les aérosols de brouillard d'huile ne peuvent être arrêtés que par des filtres submicroniques et coalescents. L'huile du compresseur est généralement piégée en sortie de production d'air comprimé. Cette huile a tendance à obstruer le réseau de distribution et est inadaptée à la lubrification des outils et équipements pneumatiques.

\subsection{Lubrification}

Le lubrificateur injecte un brouillard d’huile dans l’air comprimé. Il permet de lubrifier les pièces mobiles des équipements pneumatiques. Une huile spéciale est introduite en bout de réseau, juste avant les équipements pneumatiques.


\begin{figure}[H]
\centering
\includegraphics[width=0.5\textwidth]{Images/M21/M21_10.png}
\end{figure}



Pour améliorer la lubrification par brouillard d’huile, il faut éviter d’utiliser des flexibles trop longs entre le système de lubrification et l’équipement. En effet, le brouillard d’huile disparaît. Les fabricants recommandent une longueur maximale de 10 mètres.
 
\subsection{Les sécheurs}


\vspace{0.3cm}
\textbf{Rôle des sécheurs}
\vspace{0.3cm}


Les filtres classiques ne peuvent arrêter la vapeur d’eau. Le rôle du sécheur est donc de capter la vapeur d’eau contenue dans l’air. La quantité de vapeur à condenser est variable selon les conditions météorologiques. L’air est plus sec en hiver qu’en été.
\vspace{0.3cm}

Un air à température de 20 °C à pression atmosphérique peut contenir jusqu’à 15 grammes de vapeur d’eau ; un air à 0 °C, 3,5 grammes.


Il existe deux types de séchoirs :
\begin{itemize}

\item  les sécheurs par réfrigération, qui condensent la vapeur d’eau contenue dans l’air ;
\item  les sécheurs par adsorption, qui retiennent la vapeur d’eau grâce à un produit absorbant.

\end{itemize}


\vspace{0.3cm}
\textbf{Les sécheurs par réfrigération}
\vspace{0.3cm}

Les sécheurs par réfrigération provoquent la condensation de la vapeur d’eau. Le principe consiste à refroidir l’air humide en dessous de la température de rosée, ou température de condensation.
\vspace{0.3cm}

L’air passe sur un échangeur froid. La réfrigération est obtenue comme dans les réfrigérateurs grâce à un fluide frigorigène.


\begin{figure}[H]
\centering
\includegraphics[width=0.2\textwidth]{Images/M21/M21_11.png}
\end{figure}

\vspace{0.3cm}
\textbf{Sécheurs par adsorption}
\vspace{0.3cm } 

Les sécheurs sont constitués de deux réservoirs contenant un produit absorbant la vapeur d’eau (oxyde d’aluminium ou sels, par exemple). 

\vspace{0.3cm}

L’air comprimé est séché en traversant le réservoir. Lorsque le produit est saturé en eau, il faut le régénérer. Il existe des régénérations dites avec chaleur ou sans chaleur.
\vspace{0.3cm}


Sur le sécheur avec régénération sans chaleur, on fait circuler une faible quantité d’air sec à travers le réservoir à régénérer. Cet air, provenant du réservoir voisin, absorbe l’eau puis est évacué à l’extérieur.

\vspace{0.3cm}

Sur le sécheur à régénération par chaleur, une très faible quantité d’air sec est introduite. La régénération est en partie assurée par le chauffage du produit absorbant appelé aussi « dessiccant ».


\begin{figure}[H]
\centering
\includegraphics[width=0.4\textwidth]{Images/M21/M21_12.png}
\end{figure}




\begin{figure}[H]
\centering
\includegraphics[width=0.8\textwidth]{Images/M21/M21_13.png}
\end{figure}

\subsection{Réservoirs}

Le réservoir constitue une réserve d’air comprimé. 
Il sert :

\begin{itemize}

\item  à diminuer les fréquences de démarrage et d’arrêt des compresseurs, il évite les cycles courts et augmente ainsi la durée de vie du compresseur ;
\item  à pourvoir temporairement à une forte demande d’air comprimé ;
\item  à amortir les variations de pression dans le cas de compresseur à piston ;
\item  à séparer l’eau condensée et l’huile de l’air.

\end{itemize}

\begin{figure}[H]
\centering
\includegraphics[width=0.4\textwidth]{Images/M21/M21_14.png}
\end{figure}



Ils sont généralement équipés des composants suivants :
\begin{itemize}
\item  un manomètre permettant le contrôle de la pression de l’air stocké ;
\item  une purge automatique en point bas pour évacuer l’eau condensée ;
\item  une soupape de sécurité permettant de protéger le réseau contre les surpressions et les risques de rupture ;
\item  un trou de bras pour le nettoyage du réservoir.

\end{itemize}

\subsection{Circuit de distribution}

Les canalisations de distribution sont réalisées en métal (acier galvanisé) ou en matière plastique.



\section{Les équipements du réseau de vapeur}


\subsection{Les purgeurs d’eau}

\textbf{Nécessité de la purge}


Les purgeurs permettent d’évacuer automatiquement les condensats produits dans les canalisations et dans les corps de chauffe.
Quand la vapeur cède sa chaleur latente, elle se condense. Cette condensation se produit dans deux cas de figure.



\begin{itemize}

\item  Dans les conduites de vapeur : malgré le calorifuge, il y a des pertes thermiques qui produisent des condensats. Si les condensats ne sont pas évacués assez régulièrement, ils vont se mettre à circuler à la vitesse de la vapeur et provoquer une rapide érosion de la canalisation.
\item  Dans les échangeurs de chaleur : l’échange thermique produit des condensats. Il est alors nécessaire d’évacuer les condensats ainsi formés pour que la vapeur puisse à nouveau arriver dans l’enceinte.

\end{itemize}

\vspace{0.3cm}
\textbf{Implantation des purgeurs}
\vspace{0.3cm} 

Les purgeurs se trouvent dans deux situations principales : 
\begin{itemize}

\item  en dérivation : sur les canalisations horizontales, à tous les points bas, il est nécessaire d’avoir un point de purge précédé d’un pot de purge.
\item  en ligne : le purgeur est placé à la sortie de l’échangeur thermique, il ne doit pas y avoir de vapeur en sortie d’échangeur.

\end{itemize}

\begin{figure}[H]
\centering
\includegraphics[width=0.9\textwidth]{Images/M21/M21_15.png}
\end{figure}

\newpage
\vspace{0.3cm}
\textbf{Principes de fonctionnement des purgeurs}
\vspace{0.3cm} 

Les purgeurs travaillent de trois manières :
\begin{itemize}

\item  selon la température du fluide ;
\item  selon l’état du fluide liquide ou gazeux ;
\item  selon la vitesse du fluide.

 \end{itemize}

\vspace{0.3cm}
\textbf{Classification}
\vspace{0.3cm}

Les purgeurs peuvent être classés en trois familles :
\begin{itemize}

\item  les purgeurs thermostatiques, qui comprennent les purgeurs bimétalliques et les purgeurs thermostatiques à capsule ou à soufflets ;
\item  les purgeurs à flotteurs, qui comprennent les purgeurs à flotteurs fermés et les purgeurs à flotteurs inversés ouverts ;
\item  les purgeurs thermodynamiques. 


\end{itemize}
 
Ils peuvent également être classés en deux catégories avec des sous-catégories :




\begin{figure}[H]
\centering
\includegraphics[width=0.9\textwidth]{Images/M21/M21_16.png}
\end{figure}



\vspace{0.3 cm}
\textbf{Les purgeurs thermostatiques}
\vspace{0.3cm}

Leur principe de fonctionnement est basé sur la température du fluide dans le purgeur.
Il existe deux sous-catégories principales de purgeurs thermostatiques :


\vspace{0.3cm}
\textbf{Les purgeurs thermostatiques bimétalliques}
\vspace{0.3cm}

Dans les purgeurs thermostatiques bimétalliques l’information de température est captée par des ensembles de deux métaux ne se dilatant pas de la même façon.
\vspace{0.3cm}

Lorsque les condensats, plus froids que la vapeur, sont piégés par le purgeur, la température chute au niveau du bilame. Ce dernier se rétracte et soulève le clapet qui laisse échapper les condensats. La vapeur peu à peu remplace les condensats et réchauffe le bilame qui se dilate et obture l’orifice d’évacuation.




\begin{figure}[H]
\centering
\includegraphics[width=0.4\textwidth]{Images/M21/M21_17.png}
\end{figure}


\vspace{0.3cm}
\textbf{Les purgeurs thermostatiques à capsule ou à soufflet}
\vspace{0.3 cm}

Dans les purgeurs thermostatiques à capsule, la prise d’information de température est effectuée par une capsule remplie d’un fluide qui se dilate en fonction de la température. 

\begin{figure}[H]
\centering
\includegraphics[width=0.5\textwidth]{Images/M21/M21_18.png}
\end{figure}

\vspace{0.3cm}Pièges à
 \textbf{Les purgeurs mécaniques à flotteurs}
 \vspace{0.3 cm}
 

Le principe de fonctionnement est basé sur le niveau de liquide dans le purgeur ou la différence de densité entre la vapeur et le condensat.
Il existe deux sous-catégories principales de purgeurs à flotteurs :


\begin{figure}[H]
\centering
\includegraphics[width=0.5\textwidth]{Images/M21/M21_19.png}
\end{figure}


\vspace{0.3cm}
\textbf{Les purgeurs à flotteurs fermés}
\vspace{0.3cm}

Leur principe de fonctionnement est assez simple : le liquide dans le corps du purgeur fait flotter une boule, attelée par un bras de levier à un orifice d’obturation. Lorsque du liquide est présent dans le purgeur, le flotteur se lève et l’orifice s’ouvre permettant d’évacuer le condensat.

\vspace{0.3cm}

Ce type de purgeur permet d’évacuer le condensat dès son apparition, quelle que soit sa température. En revanche, ce purgeur ne fait pas la différence entre l’air et la vapeur. Restant fermé en présence de gaz, il peut en résulter un blocage du système sur lequel il est implanté par bouchon d’air. Dans le cas où des bouchons d’air sont à craindre, il est souhaitable de lui adjoindre un petit purgeur thermostatique.




\vspace{0.3cm}
\textbf{Les purgeurs à flotteurs inversés ouverts}
\vspace{0.3cm}

Le flotteur en forme de cloche est tourné vers le bas. Dans le fond de la cuvette, une petite quantité de condensats assure un joint d’eau. Si des condensats arrivent, le flotteur reste en position basse et l’orifice supérieur permet de les évacuer. Si de la vapeur arrive, elle se retrouve sous la cloche et la fait monter, créant un déplacement vers le haut. Le clapet lié au flotteur par une biellette assure la fermeture de l’orifice de sortie. Un petit trou est percé à la partie supérieure de la cloche pour permettre d’évacuer l’air qui pourrait éventuellement s’y trouver et engendrerait le blocage du système.



\begin{figure}[H]
\centering
\includegraphics[width=0.4\textwidth]{Images/M21/M21_20.png}
\end{figure}

\vspace{0.3cm}Pièges
\textbf{Les purgeurs mécaniques thermodynamiques}
\vspace{0.3 cm}

Le purgeur analyse la vitesse de passage du fluide. S’il s’agit de liquide, la vitesse est relativement faible ; s’il s’agit de gaz ou de vapeur, la vitesse est élevée. Cette augmentation de vitesse rapproche alors le disque de son siège et ferme le passage du fluide. 
\vspace{0.3cm}

Quand le purgeur laisse passer le fluide, les parties supérieures et inférieures du disque sont soumises à la même pression. Les surfaces prises en considération au-dessous et au-dessus du disque étant les mêmes, il n’en résulte aucun déplacement. 
\vspace{0.3cm}

Quand le purgeur se ferme, il enferme en partie supérieure et inférieure la pression de la vapeur ; cette pression s’exerçant sur des surfaces différentes, la force est plus importante sur le dessus du disque. 

\vspace{0.3cm}

C’est cette force qui permet au purgeur de rester fermé jusqu’à ce que la pression sur le dessus du disque baisse par le refroidissement de la vapeur. Le purgeur s’ouvre alors à nouveau. S’il s’agit de condensat, il laisse passer jusqu’à ce que la vapeur arrive, et il se referme s’il s’agit de vapeur.



\begin{figure}[H]
\centering
\includegraphics[width=0.4\textwidth]{Images/M21/M21_21.png}
\end{figure}


\subsection{La robinetterie manuelle de vapeur}

%\section{Vapeurs}
\vspace{0.3cm}
\textbf{Généralités}
\vspace{0.3cm}

Il existe une très grande diversité de robinetteries manuelles. Elles dépendent de la technologie, de la fréquence de l’utilisation prévue et de l’éventuelle motorisation ou de la démultiplication pour assister la fermeture de la robinetterie. 

\begin{figure}[H]
\centering
\includegraphics[width=\textwidth]{Images/M21/M21_22.png}
\end{figure}


Sauf pour la robinetterie à pointeau, la fonction première de cette robinetterie est l’isolement des équipements alimentés en vapeur.


\vspace{0.3cm}
\textbf{Le robinet à piston}
\vspace{0.3cm}

Le robinet à piston est certainement le plus utilisé dans l’industrie. Le plus connu est le robinet à piston Klinger® (marque déposée par Trouvay-Cauvin).


\begin{figure}[H]
\centering
\includegraphics[width=0.5\textwidth]{Images/M21/M21_23.png}
\end{figure}

Lorsque l’on tourne le volant de manœuvre, le mouvement de rotation est transformé en un mouvement de translation par un système de vis. Ce système de vis peut se trouver dans le chapeau ou au centre du volant. Dans ce dernier cas de figure on parle de robinet à tige sortante, la position de la tige indique alors la position de la vanne, mais le robinet possède alors une tige filetée qui dépasse du volant quand la vanne est ouverte. Cela est relativement dangereux pour le personnel qui passe à proximité. 

\vspace{0.3cm}

Le déplacement de la tige déplace un piston qui coulisse entre deux bagues d’étanchéité. Ces bagues sont réalisées grâce à un empilage de tresses graphitées et de feuillard d’acier inoxydable. Ces deux bagues sont séparées par une chemise (cylindre percé). La bague du bas assure l’étanchéité de la veine fluide, la bague du haut assure l’étanchéité vers l’extérieur de la vanne. Ces deux bagues sont tassées autour du piston par le fouloir de presse-étoupe au moyen d’un ensemble de boulons. 
\vspace{0.3cm}


Les écrous sont séparés du fouloir par des rondelles ressort permettant de reprendre les contraintes dues à la dilatation du robinet à la mise en température.
\vspace{0.3cm}

Ce robinet est assez simple à entretenir ou rénover. 


\begin{figure}[H]
\centering
\includegraphics[width=0.5\textwidth]{Images/M21/M21_24.png}
\end{figure}

\vspace{0.3cm}
\textbf{Le robinet à soupape}
\vspace{0.3cm }

Le robinet à soupape contient un opercule qui se déplace de haut en bas et assure l’étanchéité sur un siège métallique usiné. La surface de l’opercule peut être en métal usiné ou revêtue d’un joint. Ces joints sont réalisés en tresses graphitées, dans des Téflons particuliers, des caoutchoucs spéciaux ou en ensembles composites. 
\vspace{0.3cm}


L’étanchéité avec joint est en général de meilleure qualité, mais sa caractéristique se dégrade dans le temps. De plus, au-dessus de 160 °C, il existe peu de matériaux souples qui résistent. On trouve donc sur les sites principalement des robinets à étanchéité métal/métal. 
Ci-après, les différents types de robinets à soupapes.


\begin{figure}[H]
\centering
\includegraphics[width=0.9\textwidth]{Images/M21/M21_25.png}
\end{figure}


\vspace{0.3cm}
\textbf{Moyens techniques pour obtenir l’étanchéité}
\vspace{ 0.3cm}

\begin{figure}[H]
\centering
\includegraphics[width=0.4\textwidth]{Images/M21/M21_26.png}
\end{figure}

\vspace{0.3cm}
\textbf{Étanchéité vers l’extérieur}
\vspace{0.3cm}

Deux types d’étanchéité sont à distinguer dans un robinet :
\begin{itemize}

\item  les joints de bride et de couvercle (étanchéité statique) ;
\item  les joints de tige (étanchéité dynamique).

\end{itemize}

\vspace{0.3cm}
\textbf{Etanchéité statique}
\vspace{0.3cm}

Elle consiste à serrer un joint entre les faces d’appui de deux pièces fixes. Il existe quatre grandes familles de joints statiques :
\begin{itemize}

\item  le joint non métallique en feuille (élastomères comprimés, utilisation maximum à 350 °C, 40 bars) ;
\item  le joint non métallique formé (joints toriques en élastomères) ;
\item  le joint métallique (pour hautes températures et hautes pressions) ;
\item  les joints composés : métallo-plastique, téflon plastique, spiralés.

\end{itemize}

\vspace{0.3cm}
\textbf{Etanchéité dynamique}
\vspace{0.3cm}

L’étanchéité de la tige de commande de l’obturateur peut être obtenue par joint formé, par soufflet ou par presse-garniture.
Le joint formé est du type joint torique, à lobes, à lèvres, plat, etc. Ces joints sont surtout en élastomère.
Utilisation : température maximum : 150 °C. Pression < 40 bars.




\begin{figure}[H]
\centering
\includegraphics[width=0.5\textwidth]{Images/M21/M21_27.png}presse-
\end{figure}


\vspace{0.3cm}

\textbf{Étanchéité de tige avec presse-étoupe}
\vspace{0.3cm}

Cette étanchéité est assurée par un ensemble de tresses serrées sur la tige par le fouloir de presse étoupe au moyen de boulons.
\vspace{0.3cm}

Le jeu est diminué par contact direct entre la tige et la paroi.
Cette étanchéité est assez simple, mais plus on recherche une étanchéité poussée, plus les tresses seront serrées sur la tige et plus la manœuvre du volant sera dure. Cette étanchéité sera amoindrie à chaque manœuvre et il sera nécessaire de reprendre le serrage du fouloir de presse-étoupe assez régulièrement.



\vspace{0.3cm}
\textbf{Étanchéité de tige avec soufflet}
\vspace{0.3cm}

Pour éviter les inconvénients, d’entretien d’un presse-étoupe classique et la difficulté à tourner le volant de manœuvre, la plupart des fournisseurs proposent de la robinetterie où l’étanchéité vers l’extérieur est réalisée par un soufflet métallique ou plastique, dont une extrémité est soudée sur l’opercule et l’autre serrée sur le corps du robinet.
\vspace{0.3cm}

En cas de rupture du soufflet, un presse-étoupe de sécurité traditionnel est quand même prévue. Sa fonction étant uniquement de limiter les dégâts, il n’a donc pas à être très étanche et ne durcit pas la manœuvre du volant. Ce robinet est plus cher que 
celui avec une étanchéité traditionnelle.




\vspace{0.3cm}
\textbf{Robinets à tournant cylindrique ou conique}
\vspace{0.3cm}

On distingue :
\begin{itemize}
\item  les robinets à tournant conique à boisseau foncé ;
\item  les robinets à tournant conique à boisseau défoncé ;
\item  les robinets à tournant conique à boisseau inversé ;
\item  les robinets à tournant sphérique libre (applications courantes et petits diamètres) ;
\item  les robinets à tournant sphérique guidés (applications hautes températures et grands diamètres).

 \end{itemize}

Une bonne étanchéité est obtenue par un ajustement soigné entre le tournant et le corps.


Avantages 
\begin{itemize}

\item  commande rapide (un quart de tour) ;
\item  faible perte de charge ;
\item  peu d’entretien ;
\item  encombrement réduit ;
\item  très bonne étanchéité.

\end{itemize}

Inconvénients 
\begin{itemize}

\item  les robinets à boisseau ne permettent pas le réglage du débit ;
\item  la température de service est limitée à 180 °C pour les fournitures courantes ;
\item  complexité des robinets de grand diamètre.

\end{itemize}

\vspace{0.3cm}
\textbf{Robinets à papillon}
\vspace{0.3cm}

Ces robinets ont pris une place importante en robinetterie industrielle. On distingue les robinets :
\begin{itemize}
\item  à papillon centré ou excentré ;
\item  à joint sur le papillon, sur le corps, ou à manchette.

\end{itemize}


\begin{figure}[H]
\centering
\includegraphics[width=0.35\textwidth]{Images/M21/M21_28.png}
\end{figure}

\newpage
\textbf{Avantages }
\begin{itemize}
\item  faible encombrement ;
\item  simplicité liée au peu de pièces constitutives ;
\item  prix de revient faible ;
\item  très bonne étanchéité, notamment avec les joints en élastomère ;
\item  faible perte de charge ;
\item  aptitude au réglage et aux manœuvres fréquentes.


\end{itemize}

Inconvénients 
\begin{itemize}

\item  pas de passage intégral ;
\item  pas adapté au fonctionnement à forte pression différentielle. 

\end{itemize}

\vspace{0.3cm}

\textbf{Robinets vannes}
\vspace{0.3cm}

On distingue deux familles : les robinets à vannes à sièges obliques et ceux à vannes à sièges parallèles.


\begin{figure}[H]
\centering
\includegraphics[width=0.4\textwidth]{Images/M21/M21_29.png}
\end{figure}


\textbf{Vannes à sièges obliques}
\vspace{0.3cm}

L’étanchéité est obtenue par coincement de l’opercule entre les portées d’étanchéité obliques du corps.


\newpage
\textbf{Vannes à opercule monobloc}
\vspace{0.3cm}

\textbf{Avantages }
\begin{itemize}

\item  construction est simple ;
\item  grande précision d’usinage nécessaire pour obtenir une bonne étanchéité.

\item \end{itemize}

\textbf{Inconvénients }
\begin{itemize}

\item  l’obturateur massif ne convient pas aux utilisations à haute température ;
\item  la forte pression peut provoquer des déformations du corps dans les grands diamètres.

\end{itemize}

\vspace{0.3cm}
\textbf{Vanne à opercule élastique}
\vspace{0.3cm}

Meilleure tenue en température élevée et en grand diamètre.

\vspace{0.3cm}
\textbf{Vanne à double opercule} 
\vspace{0.3cm}

L’obturateur est constitué de deux éléments appuyés l’un contre l’autre. Cette disposition permet de compenser la variation angulaire des portées d’étanchéité.
Le domaine d’utilisation est plus étendu que celui des vannes à opercule élastique.


\vspace{0.3cm}
\textbf{Vanne à sièges parallèles}
\vspace{0.3cm}

Les disques de l’opercule frottent pendant la totalité de la course sur les sièges.
Attention au choix des matériaux utilisés pour éviter une usure prématurée.



Avantages 
\begin{itemize}

\item  peu de sensibilité aux déformateurs thermiques ou dues à la pression ;
\item  facilité d’entretien

\end{itemize}

\vspace{0.3cm}
\textbf{Vanne à opercule à libre dilatation}
\vspace{0.3cm}

Action d’un ressort intérieur qui applique les deux disques de l’opercule sur leur siège.

\begin{figure}[H]
\centering
\includegraphics[width=0.25\textwidth]{Images/M21/M21_30.png}
\end{figure}

\vspace{0.3cm}
\textbf{Vanne à opercule à dispositif de blocage}
\vspace{0.3cm}

En position de fermeture, les deux disques de l’obturateur sont bloqués sur leur siège.


Avantages des robinets-vannes
\begin{itemize}
\item  robustesse ;
\item  domaine étendu d’application en température et pression (portées d’étanchéité métallique) ;
\item  adaptés à l’isolement des circuits ;
\item  faible perte de charge.

\end{itemize}

Inconvénients des robinets-vannes
\begin{itemize}

\item  encombrement et poids importants ;
\item  ne conviennent pas au réglage, aux fluides chargés ;
\item  non adaptés aux manœuvres fréquentes ;
\item  étanchéité poussée difficile à obtenir. 
 \end{itemize}


Ce type de robinet offre une étanchéité assez médiocre et reste assez volumineux. Il est assez donc peu utilisé sur les installations de vapeur.
\vspace{0.3cm}

Ce type de robinet à un opercule de forme conique, l’opercule est en général en acier traité en surface et sa fonction première est de créer un réglage de débit de fluide 


\subsection{Vannes de régulation}

Les vannes de régulation permettent de régler automatiquement le débit de vapeur de la chaudière et de ses équipements.


\vspace{0.3cm}
\textbf{Caractéristique des clapets}
\vspace{0.3cm}


Selon le type de clapet, une même levée ne laisse pas passer le même débit de vapeur. Il existe trois types de vannes de régulation : à ouverture rapide, à ouverture linéaire, à égal pourcentage.


\begin{figure}[H]
\centering
\includegraphics[width=0.3\textwidth]{Images/M21/M21_31.png}
\end{figure}

\vspace{0.3cm}
\textbf{Ouverture rapide}
\vspace{0.3cm }

Le clapet est plat. Dès que le clapet se soulève de son siège, le débit est important. Ce robinet ne permet pas un réglage précis du débit.

\vspace{0.3cm}

\textbf{À ouverture linéaire et à égal pourcentage}
\vspace{0.3cm }

La vanne permet un réglage précis du débit fonction de la levée du clapet. Le choix de la vanne se fera en fonction de la caractéristique de l’échangeur.


\vspace{0,3cm}
\textbf{Sens de fonctionnement des moteurs}
\vspace{0,3cm}


Selon le type de vanne, une montée de la tige donnera une augmentation de débit pour une vanne à mono-siège direct, une diminution de débit pour une vanne à mono-siège inverse.
\vspace{0.3cm}

Les vannes de type mono-siége inverse sont parfois appelées autoclaves, ce qui indique que ces vannes se fermeront d’elles-mêmes, sans moteur, par le seul effet de la pression du fluide.


\begin{figure}[H]
\centering
\includegraphics[width=0.9\textwidth]{Images/M21/M21_15.png}
\end{figure}


\vspace{0.3cm}
\textbf{Technologie de vanne de régulation}
\vspace{0.3 cm}

Les différentes remarques présentées dans le paragraphe sur les vannes manuelles sont pour la plupart transposables ici. Les particularités des vannes de régulation sont la forme du clapet plus élaborée, en général au détriment des caractéristiques d’étanchéité. 
Voici ci-dessus les principales formes de clapets présents sur le marché.
\vspace{0.3cm}

Le presse-étoupe (étanchéité de la tige) est en général plus élaboré car la vanne de régulation manœuvrera beaucoup plus souvent que la vanne manuelle. Il est donc nécessaire d’avoir une étanchéité beaucoup plus résistante à l’érosion.



\section{Les tuyauteries et raccords}

\subsection{L'acier}


L’acier est le matériau le plus utilisé dans les transports de fluide. Il existe différents modes de fabrication des canalisations : par soudage ou par étirage.
Le tube obtenu par soudage est un assemblage de plaques roulées et soudées. Ce type de canalisation ne résiste pas aux fortes pressions nécessaires au fonctionnement de la chaudière vapeur.
\vspace{0.3cm}

Pour les pressions élevées, on utilisera plutôt du tube acier étiré obtenu à partir d’un lingot d’acier chauffé à blanc.
\vspace{0.3cm}

La résistance de la tuyauterie aux fortes pressions dépend de son épaisseur : plus l’épaisseur de la canalisation est élevée, plus elle résiste à la pression du fluide.
\vspace{0.3cm}

L’acier est un métal ferreux. On notera qu’il existe aussi pour les transports de fluide des métaux non ferreux, tels que le cuivre, l’inox ou le nickel, qui résistent particulièrement bien à la corrosion, mais qui, en raison de leur coût élevé, ne sont pas employés dans les usines.
La dénomination des canalisations acier est normalisée et facilite leur choix.


\subsection{Les méthodes de raccordement des tuyaux
}

Il existe trois méthodes pour raccorder les canalisations :
\begin{itemize}
\item  le filetage et les raccords filetés ;
\item  le soudage ;
\item  l’emploi de brides.

\end{itemize}

\vspace{0.3cm}
\textbf{Le filetage}
\vspace{0.3cm}

À chaque extrémité de la canalisation, on réalise un filetage qui permet un montage facile à l’aide de raccords filetés. Le filetage est utilisé lorsque les pressions dans le réseau ne sont pas trop élevées, car la présence des filets réduit la résistance de la tuyauterie.
Les tuyaux sont coupés à l’aide d’un coupe-tube et le filetage est réalisé à l’aide d’une 
filière.


\vspace{0.3cm}
\textbf{Raccords filetés}
\vspace{0.3cm}

Il existe un grand choix de raccords filetés pour l’assemblage des canalisations :
\begin{itemize}
\item  les coudes, qui permettent un changement de direction ;
\item  les unions, qui assurent un montage et démontage rapide ;
\item  les manchons, qui permettent la jonction de deux canalisations droites de grande longueur ;
\item  les tés et les croix, qui permettent de réaliser des embranchements ;
\item  les mamelons, qui permettent la jonction de deux canalisations ;
\item  les réductions, qui permettent de raccorder deux canalisations de diamètres différents ;
\item  les bouchons, qui permettent d’obturer une canalisation en attente.

\end{itemize}

\vspace{0.3cm}
\textbf{Soudage}
\vspace{0.3cm} 

Le soudage permet de raccorder les tuyauteries bout à bout et aussi de les assembler aux raccords ou aux équipements tels que les vannes ou les robinets.
Par rapport à l’assemblage par filetage ou par bride, le soudage permet :

\begin{itemize}

\item  d’éviter les fuites liées à la détérioration des joints ;
\item  de faciliter l’isolation thermique des canalisations.

 \end{itemize}

Les raccords et les accessoires à souder possèdent un évidement dans lequel s’insère le tuyau lisse. Différentes méthodes de soudage sont employées : à l’arc, brasure autogène, hétérogène…


\vspace{0.3cm}
\textbf{Les brides}
\vspace{0.3cm} 

Les brides sont montées à chaque extrémité de la canalisation, puis assemblées à l’aide de boulons. Ce mode d’assemblage est beaucoup plus résistant que le filetage. L’étanchéité est meilleure qu’un joint fileté dans la mesure où l’alignement des canalisations est parfaitement réalisé.
\vspace{0.3cm}

Il existe plusieurs types de brides :

\begin{itemize}
\item  les brides filetées qui sont vissées à l’extrémité de la tuyauterie ;
\item  les brides soudées sur la canalisation ;
\item  les brides coulissantes à bords surélevés.

\end{itemize}

\subsection{La dilatation}

Lorsque les matériaux s’échauffent, ils se dilatent ou s’allongent. Rien ne résiste à la dilatation des matériaux et pour absorber le phénomène de dilatation, on a recours à des lyres de dilatation ou à des joints d’expansion.
Dans le cas des lyres, l’allongement est absorbé par la flexion des bras de la lyre.

\vspace{0.3cm}

Dans le cas des joints d’expansion, deux parties emboîtées l’une dans l’autre coulissent. Le joint d’expansion nécessite l’emploi de lubrifiant et de garnitures d’étanchéité.


\subsection{Supports}

L’ensemble de la tuyauterie doit être supporté afin que sa masse ne repose pas sur 
les équipements de l’usine. Les supports sont judicieusement placés pour éviter tout fléchissement.

\vspace{0.3cm}

En matière de supports, il existe les colliers de suspension, les supports et consoles de 
fixation.


\section{Les installations secourues}


\subsection{Généralités}


Certaines installations ne peuvent être mises à l’arrêt en cas de défaillance du réseau de distribution EDF (pannes ou autre), soit pour des raisons de sécurité humaine, soit pour des raisons économiques. La solution consiste alors à assurer la production d’énergie électrique de façon autonome, en se substituant au réseau EDF.
\vspace{0.3cm}


Pour satisfaire cette contrainte, il existe plusieurs types d’appareillage. Les plus utilisées sont :

\begin{itemize}

\item  le groupe de secours ou groupe électrogène (GE) ;
\item  l’alimentation sans interruption (ASI) ou « onduleur ».

\end{itemize}


\subsection{Le groupe électrogène (GE)}

\begin{figure}[H]
\centering
\includegraphics[width=0.9\textwidth]{Images/M21/M21_33.png}
\end{figure}


\vspace{0.3cm}
\textbf{Les différents éléments constitutifs}
\vspace{0.3cm}

Cet équipement est constitué par :
\begin{itemize}

\item  Un moteur Diesel de puissance adaptée aux besoins de l’installation, équipé de circuits auxiliaires :
    \begin{itemize}
    \item  un circuit de démarrage constitué d’une chaîne de démarrage comprenant un démarreur, une batterie et son chargeur.
    \item  un circuit de gas-oil comprenant : 

        \begin{itemize}
        \item  une cuve dite journalière de capacité maximale d’environ 500 l ;
    	\item  une cuve extérieure de capacité déterminée en fonction de l’autonomie demandée au moteur ;
    	\item  une pompe de gas-oil automatique doublée d’une pompe manuelle assurant le remplissage de la cuve journalière à partir de la cuve extérieure.
        \end{itemize}
        
    \item  un circuit de graissage équipé d’une réserve d’huile déterminée en fonction de l’autonomie demandée au moteur ;
    \item  un circuit de refroidissement, à air ou à eau selon le type de moteur. Lors d’un refroidissement à air, c’est l’arbre moteur qui entraîne directement ou par courroies l’arbre du ventilateur. Lors d’un refroidissement à eau, il y a nécessairement un échangeur, une pompe de circulation et un ventilateur.

    \end{itemize}



\item  Un alternateur de puissance adapté aux besoins de l’installation, équipé d’un système de régulation de tension et de fréquence.
\end{itemize}



Au niveau de la liaison électrique, on retrouve le plus souvent les équipements suivants : 
\begin{itemize}

\item  un transformateur élévateur, qui rehausse la tension de 400 à 5 ou 20 kV ;
\item  un disjoncteur assurant le découplage du réseau EDF et le couplage du GE sur l’installation ; on le nomme disjoncteur « Normal-Secours » (NS) ;
\item  un système de délestage permettant d’alimenter en priorité les réseaux sensibles et assurant l’ouverture des disjoncteurs des autres réseaux.


\end{itemize}

\vspace{0.3cm}
\textbf{Fonctionnement}
\vspace{0.3cm}

Après une détection d’absence de tension sur le secteur (20 secondes), le groupe électrogène démarre (50 secondes), le réseau de délestage se met en configuration et le disjoncteur Normal-Secours permute, le moteur prend la charge plus ou moins progressivement par la fermeture des disjoncteurs prioritaires (relestage, 210 secondes).
\vspace{0.3cm}

La mise en rotation de l’arbre moteur, permet à l’alternateur d’être la source de phénomènes magnétiques assurant la transformation de l’énergie mécanique du moteur en énergie électrique sous forme de tensions alternatives triphasées.


\vspace{0.3cm}
\textbf{Conclusion}
\vspace{0.3cm}


Cet équipement est très utilisé dès que la puissance de secours est importante ou que le risque de coupure est fréquent. 
\vspace{0.3cm}

En le dimensionnant de façon précise, on peut satisfaire les besoins de toute installation ; en revanche, on ne peut assurer un secours sans interruption. En effet, de par son principe de fonctionnement, il faut qu’il y ait absence de secteur pour que le groupe électrogène démarre.
Pour répondre à cette contrainte supplémentaire, on utilise les alimentations sans interruption, notées ASI, appelées aussi « onduleurs ».



\subsection{L’alimentation sans interruption (ASI)}

\textbf{Utilité}
\vspace{0.3cm}

Cet équipement est utilisé dans le but d’assurer un secours en cas de coupure du réseau distributeur, mais il est aussi préconisé afin d’obtenir un « courant propre », c’est-à-dire sans trop d’harmoniques (cf. module 17). Il est pratiquement indispensable pour l’alimentation des salles informatiques ou de gros ordinateurs.

\newpage
\vspace{0.3cm}
\textbf{Les différents éléments constitutifs}
\vspace{0.3cm}

Cet équipement est constitué de trois éléments principaux :
\begin{itemize}

\item  un redresseur chargeur, dont le rôle est de transformer la tension alternative délivrée par le réseau EDF en une tension continue permettant la charge des batteries ;
\item  des batteries d’accumulateurs, qui assureront l’énergie électrique lors de la coupure du secteur ; l’autonomie de l’ASI dépendra uniquement de la capacité de stockage des accumulateurs ;
\item  un onduleur qui retransforme la tension continue en tensions alternatives triphasées.

\end{itemize}



\begin{figure}[H]
\centering
\includegraphics[width=\textwidth]{Images/M21/M21_34.png}
\end{figure}


\textbf{Fonctionnement}
\vspace{0.3cm}

Lorsque le secteur est présent, il assure la charge des accumulateurs. Le passage au travers des différents éléments constitutifs de l’ASI assure une alimentation « saine » et stable du réseau situé en aval.
\vspace{0.3cm}

Lorsque le secteur vient à manquer, les batteries d’accumulateurs prennent le relais de façon instantanée, sans coupure, et assurent une alimentation « saine » et stable du réseau situé uniquement en aval.
\vspace{0.3cm}

Bien entendu, si l’absence du secteur vient à durer, les batteries se déchargent et l’alimentation du réseau situé en aval n’est plus garantie.


\vspace{0.3cm}
\textbf{Conclusion}
\vspace{0.3cm}

Pris séparément, GE et ASI ont des avantages complémentaires. C’est pourquoi, sur les installations secourues, on retrouve le plus souvent ces deux équipements.
Les ASI permettent l’alimentation sans interruption des réseaux sensibles, le GE permet l’alimentation des autres réseaux plus puissants, tout en assurant la relève des batteries d’accumulateurs.
\vspace{0.3cm}

L’opération de maintenance et d’entretien primordiale pour ces équipements consiste à vérifier périodiquement que les batteries d’accumulateurs sont bien en charge, aussi bien sur les ASI que sur le ou les GE.



\section{Les moyens de lutte contre l’incendie}


\subsection{Introduction}

Les moyens de défense contre l’incendie ne s’improvisent pas. Pour combattre le feu avec le minimum de dégâts, il importe d’agir vite, ce qui implique :
\begin{itemize}

\item  un personnel parfaitement instruit des différentes mesures de protection à prendre et de l’emploi des divers moyens de lutte contre le feu ;
\item  des moyens de détection, de mise en sécurité (alarme) et d’alerte ;
\item  un matériel d’extinction approprié et toujours disponible.

\end{itemize}

« Le feu s’éteint dans la première minute avec un verre d’eau, 
dans la deuxième minute avec un seau d’eau, 
dans la troisième minute avec une tonne d’eau, 
après… »


\subsection{Les moyens de détection du feu}


\textbf{Les principaux types de détecteurs}
\vspace{0.3cm}


Le développement d’un feu se caractérise par différentes phases :
\begin{itemize}

\item  des odeurs de brûlé ou de roussi (dues à des gaz de combustion invisibles) ;
\item  des fumées plus ou moins sombres, composées de minuscules particules en suspension dans l’air ;
\item  des flammes (radiations) ;
\item  une élévation sensible de la température.

\end{itemize}

Selon le type de feu, ces quatre phénomènes se produisent plus ou moins tôt et de façon plus ou moins importante. C’est pourquoi une bonne évaluation des risques permet de choisir le détecteur le mieux adapté aux feux susceptibles de se déclarer dans les divers secteurs d’un bâtiment. Pour cela, on distingue trois grands types de détecteurs :

\begin{itemize}
\item  les détecteurs de fumées ;
\item  les détecteurs de chaleur ;
\item  les détecteurs de flammes.
\end{itemize}

À ce jour, seulement un faible pourcentage de nos UIOM est équipée de détecteurs. Ces détecteurs se trouvent dans des endroits bien précis : hall de déchargement-fosse, local compresseur, local turbo, zones de stockage en tous genres, trappes de désenfumage ou pyrodômes.

\vspace{0.3cm}
\textbf{Détecteurs de fumée}
\vspace{0.3cm}

Le détecteur ionique
\vspace{0.3cm}

Sensible à tous les types d’aérosols (particules invisibles, fumées claires ou sombres…), il agit comme un véritable « nez électronique ». Jusqu’à ces dernières années, c’était le détecteur le plus couramment utilisé (environ 80\% des applications). Il s’avère très efficace dans le cas de feux couvants à développement lent. Il fonctionne selon le principe de la chambre d’ionisation et il utilise l’influence perturbatrice générée par les produits de combustion sur de l’air ionisé, air rendu conducteur par une pastille radioactive (suffisamment faible pour ne pas causer de souci en matière de radioprotection). Les produits de combustion ont pour effet de ralentir le mouvement des ions. L’intensité du courant diminue et, à partir d’un certain seuil, le détecteur déclenche l’alarme.

Le détecteur optique à diffusion de lumière (effet Tyndall)
\vspace{0.3cm}

Son principe de fonctionnement consiste dans la visualisation d’un faisceau de lumière par la charge d’aérosols qu’il traverse (fumées ou particules). Lorsqu’un rayon lumineux rencontre un aérosol, la lumière est diffusée dans toutes les directions et captée par une cellule. Cette cellule génère un signal électrique qui déclenche l’alarme. 
Ce détecteur ne convient pas aux atmosphères poussiéreuses (hall de réception, zone de tri…).

Le détecteur optique linéaire (à absorption ou à atténuation)
\vspace{0.3cm}

Ce matériel comporte un émetteur de faisceau infrarouge et un récepteur. Lorsque le faisceau rencontre des aérosols, le rayonnement reçu est atténué. Si celui-ci s’affaiblit au-dessous d’une valeur déterminée, l’alarme se déclenche. 
\vspace{0.3cm}

On utilise surtout ce type de détecteur pour les grands volumes, lorsque le sol est encombré ou si l’installation d’un détecteur ponctuel est difficile dans la configuration du lieu.
Dans nos métiers, les détecteurs optiques, particulièrement sensibles aux poussières, provoquent très fréquemment de fausses alarmes.

\vspace{0.3cm}
\textbf{Détecteurs de chaleur}
\vspace{0.3cm}prédéfinie

Le détecteur thermostatique
\vspace{0.3cm}

Les détecteurs thermostatiques déclenchent l’alarme lorsque la température ambiante atteint une valeur fixée à l’avance, souvent lorsque le feu est déjà très important. 
Ce type de détecteur peut être utilisé sur les fosses à ordures ménagères et sur certains centres de compostage.

Le détecteur thermovélocimétrique
\vspace{0.3cm}

Ces détecteurs de chaleur réagissent en fonction de la vitesse d’élévation de la température, d’où leur nom de détecteurs thermovélocimétriques. Si la température varie fortement en peu de temps, l’alarme est donnée. 
Ce type de détecteur est peu ou pas utilisé dans nos UIOM.


\vspace{0.3cm}
\textbf{Détecteurs de flammes}
\vspace{0.3cm} 


Le détecteur optique à infrarouge ou ultraviolet
\vspace{0.3cm} 

Ces détecteurs réagissent au rayonnement modulé émis par les flammes, sans tenir compte des nombreux rayons parasites dans l’environnement de l’entreprise (soleil, appareils de chauffage, lampes à incandescence, arcs électriques…). Ils mettent en œuvre un traitement du signal très élaboré. 
\vspace{0.3cm}

Les détecteurs optiques de flammes sont généralement utilisés dans les locaux de grand volume et sont adaptés à la surveillance des fosses de réception des déchets.


\vspace{0.3cm}
\textbf{Les évolutions}
\vspace{0.3cm} 

Les détecteurs deviennent multicritères. Plusieurs capteurs ont été intégrés dans ces détecteurs conventionnels, les signaux de ces capteurs sont combinés pour en déduire l’information « feu » selon des règles définies. Ces systèmes sont conçus pour limiter les alarmes injustifiées.


\subsection{Matériel de première intervention}

Dans l’entreprise, les moyens de lutte contre le feu lors de la première intervention sont principalement des extincteurs mobiles (portatifs et sur roues) et des robinets d’incendie armés.

\vspace{0.3cm}
\textbf{Extincteurs}
\vspace{0.3cm}


Moyen de première intervention par excellence, l’extincteur portatif a souvent un rôle décisif dans la lutte contre l’incendie. Chaque année, il vient à bout d’un nombre considérable de foyers naissants qui, en son absence, auraient dégénéré en sinistres. C’est pourquoi la rapidité d’action est primordiale 
Dans d’autres cas, il permet de limiter les dégâts dans l’attente de la mise en œuvre de moyens plus puissants.
\vspace{0.3cm}

Un extincteur est un appareil qui permet de projeter et de diriger sur un foyer d’incendie un agent extincteur. La projection hors de l’appareil est obtenue par l’effet d’une pression intérieure qui peut être due à la compression préalable de l’agent extincteur ou à la libération d’un gaz de chasse au moment de la mise en œuvre. L’attaque du feu s’effectue toujours à sa base.
Pour attaquer efficacement un début d’incendie, il faut disposer de l’agent extincteur le mieux approprié à la nature du feu.
\vspace{0.3cm}

La norme AFNOR NF EN 2 distingue quatre classes de feu (voir tableau) :
A. Feux secs de solides,
B. Feux gras de liquides ou solides liquéfiables
C. Feux de gaz
D. Feux de métaux

\vspace{0.3cm}


L’extincteur est essentiellement caractérisé en fonction de l’agent extincteur qu’il contient. À l’heure actuelle, on distingue :
\begin{itemize}
\item  les extincteurs à base d’eau (pulvérisée, avec ou sans additif, et sous forme de mousse) ;
\item  les extincteurs à poudres ABC ;
\item  les extincteurs à dioxyde de carbone (C02),
\item  les extincteurs à hydrocarbures halogénés (uniquement en postes fixes).

\end{itemize}

\vspace{0.3cm}
\textbf{Agents extincteurs}
\vspace{0.3cm}

\textbf{Eau}
\vspace{0.3cm} 

L’eau est le plus utilisé des agents extincteurs car on peut toujours, sauf cas exceptionnel, s’en procurer. Elle agit doublement :

\begin{itemize}

\item  par action directe, en étouffant le foyer (à l’aide de l’eau et de la vapeur formée) ;
\item  par action indirecte, en refroidissant les matériaux en combustion et en limitant les effets thermiques de l’incendie.

\end{itemize}

L’eau pulvérisée augmente considérablement l’effet de refroidissement par une vaporisation plus intense et diminue l’effet du rayonnement.
\vspace{0.3cm}

Projetée au moyen de lance, en « jet plein » (ou « jet bâton »), l’eau convient bien aux feux de classe A. Elle produit un effet mécanique qui favorise la pénétration du foyer et la dispersion des matériaux. L’utilisation du « jet plein » est à déconseiller sur les installations électriques.

Attention !
\vspace{0.3cm}

L’eau est strictement prohibée comme moyen d’extinction dans certains types d’industries (fonderies d’aluminium, par exemple), en particulier lorsque sa vaporisation trop rapide entraîne la projection de matériaux liquides (métaux, sels fondus…).


\vspace{0.3cm}
\textbf{Eau et additifs}
\vspace{0.3cm}


Pour accroître le pouvoir extincteur de l’eau, on peut lui adjoindre des tensioactifs (composé chimique susceptible d’augmenter les propriétés d’étalement, de mouillage d’un liquide).
L’eau et ces additifs se rencontrent principalement dans les extincteurs portatifs pour agir sur les feux des classes A et B.


\vspace{0.3cm}
\textbf{Les poudres}
\vspace{0.3cm} 

La projection de la poudre est obtenue par la pression d’un gaz auxiliaire.
Les appareils à poudre agissant principalement par étouffement et/ou inhibition sont plus efficaces en local clos qu’en plein air. Bien qu’ils puissent être utilisés en présence de courant électrique, il est déconseillé de les mettre en œuvre directement sur le matériel électrique (risques de détérioration).
La majorité des poudres ne présentent qu’un faible risque toxicologique pour l’homme ; toutefois, elles sont en général légèrement irritantes pour les voies respiratoires et les muqueuses. La projection de la poudre dans un local diminue fortement la visibilité.


\vspace{0.3cm}
\textbf{Gaz inertes}
\vspace{0.3cm} 

L’extinction avec les gaz inertes (dioxyde de carbone, azote, leurs mélanges…) est obtenue par diminution de la teneur en oxygène dans l’atmosphère.

Quelle que soit la nature de l’agent extincteur qu’ils contiennent, les appareils sont répartis en deux grandes catégories, en fonction de leur masse en ordre de marche :


\begin{itemize}
\item  les extincteurs portatifs, dont la masse totale est inférieure ou égale à 20 kg, utilisables par toute personne en présence d’un début d’incendie ;
\item  les extincteurs sur roues, qui peuvent être tractables à bras, remorquables ou à la fois tractables à bras et remorquables. Ce sont des appareils de première intervention susceptibles d’être mis en œuvre par un service de sécurité incendie.

 \end{itemize}
 
Du point de vue quantitatif, la charge des appareils à base d’eau s’exprime en volume (litres) et celle des autres appareils en masse (kg).

Les extincteurs sont également caractérisés par :

\begin{itemize}

\item  leur efficacité, c’est-à-dire leur aptitude à éteindre, dans une ou plusieurs classes de feux déterminées, un foyer d’une importance définie ;
\item  leur durée totale de fonctionnement, c’est-à-dire le temps pendant lequel a lieu la projection de l’agent extincteur, sans qu’il y ait eu d’interruption dans la projection, la vanne étant totalement ouverte et l’émission du gaz de chasse n’étant pas prise en considération.


\end{itemize}


Un extincteur doit être impérativement adapté à la nature des combustibles et aux risques encourus. 
La norme NF EN 2 distingue quatre classes de feux :




\begin{figure}[H]
\centering
\includegraphics[width=\textwidth]{Images/M21/M21_35.png}
\end{figure}

	

\vspace{0.3cm}
\textbf{Les robinets d’incendie armés (RIA)}
\vspace{0.3cm} 

Les postes d’incendie, généralement désignés sous le sigle RIA (robinets d’incendie armés), sont des installations fixes de premier secours contre l’incendie. Destinés à être mis en œuvre par une seule personne dès l’alerte, ils sont au tout premier rang des matériels d’autoprotection.

Le robinet d’incendie armé se compose :

\begin{itemize}

\item  d’un dévidoir à alimentation axiale ;
\item  d’une longueur élémentaire de 20 ou 30 m de tuyau semi-rigide.

\end{itemize}

Le réseau d’eau public ne peut être utilisé que si le débit requis par l’installation n’excède pas 80\% du débit disponible de la canalisation d’eau de ville aux heures de plus grande consommation. De plus, le diamètre nominal de cette canalisation doit être supérieur ou égal au DN 80.
\vspace{0.3cm}

Le niveau de l’eau des réservoirs et bacs doit pouvoir être vérifié par un indicateur visible de l’extérieur et l’insuffisance d’eau signalée par un appareil donnant une alarme audible. Réservoirs et bacs doivent être protégés contre le gel.
\vspace{0.3cm}

La pression minimale au robinet diffuseur du RIA le plus éloigné doit être de 0,25 MPa (2,5 bars), avec généralement un maximum de quatre robinets en fonctionnement simultané.
\vspace{0.3cm}

Les robinets d’incendie armés permettent, lorsque l’emploi de l’eau n’est pas interdit, une action souvent puissante et efficace lors de la première intervention, dans l’attente des moyens plus importants. Ils font partie des installations fixes. 
\vspace{0.3cm}

Ils doivent :

\begin{itemize}
\item  être implantés de telle sorte que chaque point de la surface à protéger puisse être atteint par deux jets au moins ;
\item  pouvoir être alimentés en eau avec une pression minimale de 2,5 bars au robinet le plus élevé ou le plus défavorisé.

\end{itemize} 

Les RIA doivent être implantés à des emplacements abrités du gel et à proximité des accès. Ils sont signalés de façon claire.

\vspace{0,3cm}
\textbf{Les colonnes sèches ou en charge (dites humides)}
\vspace{0,3cm}

Les colonnes sèches et les colonnes en charge (dites colonnes humides) sont des tuyauteries fixes et rigides installées à demeure dans les constructions et destinées à être raccordées aux tuyaux des sapeurs-pompiers.
Les colonnes sèches sont mises en charge au moment de l’emploi.
Les colonnes en charge (dites colonnes humides) sont reliées à des réservoirs et à des pompes, à des surpresseurs ou à tout autre dispositif permettant d’alimenter les lances des sapeurs-pompiers.



\vspace{0.3cm}
\textbf{Les colonnes humides ou en charge}
\vspace{0.3cm}


Une installation de protection incendie par colonne en charge comprend au moins :

\begin{itemize}
\item  une réserve d’eau alimentée en eau (au moins 120 m3 d’eau, répartis en plusieurs cuves, exclusivement réservés aux services d’incendie) ;
\item  un dispositif d’alimentation (débit de 1 000 l/min par colonne, deux colonnes fonctionnant simultanément, pression minimale par colonne entre 4,5 et 8,5 bars, une alimentation automatique par une ou plusieurs sources de courant électrique autonomes et indépendantes de celles utilisées en service normal) ;
\item  une colonne en charge par escalier.

\end{itemize}

Ces installations permettent d’accélérer le temps d’intervention des services de secours en évitant des déploiements verticaux de tuyaux souples.

\vspace{0.3cm}
\textbf{Les colonnes sèches}
\vspace{0.3cm}

Une colonne sèche comprend :

\begin{itemize}

\item  un raccord d’alimentation ;
\item  éventuellement un élément de conduite reliant le raccord d’alimentation à la colonne proprement dite ;
\item  la colonne proprement dite ;
\item  des prises d’incendie simples ou doubles.

\end{itemize}

\vspace{ 0.3cm}
\textbf{Sprinklers}
\vspace{0.3cm}

Le rôle d’une installation de sprinklers est de déceler un foyer d’incendie, de donner une alarme et de l’éteindre à ses débuts ou au moins de le contenir de façon que l’extinction puisse être menée à bien par les moyens de l’établissement protégé ou par les sapeurs-pompiers.
\vspace{0.3cm}

La détection est faite par les sprinklers, des détecteurs thermiques à température fixe. Le choix des températures et des emplacements est fonction des caractéristiques des locaux et de leur utilisation.
\vspace{0.3cm}

L’alarme est donnée par l’intermédiaire d’une turbine hydraulique actionnant le gong d’alarme et par un pressostat de report d’alarme électrique placé à chaque poste de contrôle. 
L’extinction est effectuée par l’eau déversée par les sprinklers. 
\vspace{0.3cm}

Le débit d’eau, le type de pulvérisation, la surface d’arrosage de chaque sprinkler sont fonction du feu à éteindre. L’efficacité maximum est obtenue quand tout incendie peut être  combattu par arrosage à ses débuts, quel que soit son emplacement et avec la quantité d’eau nécessaire.


\vspace{0.3cm}
\textbf{Systèmes spéciaux}
\vspace{0.3cm}

\textbf{Canons (à eau ou à mousse)}
\vspace{0.3cm} 

L’installation de canons à eau permet de protéger efficacement les fosses de réception des ordures ménagères et les trémies d’enfournement de nos incinérateurs. 
L’eau pulvérisée est toujours accompagnée d’un additif. Le pourcentage d’additif varie selon l’effet recherché :

\begin{itemize}
\item  1\% d’additif pour la protection des fosses (pouvoir mouillant = baisse de la température) ;
\item  3\% d’additif pour la protection des trémies (mousse = pouvoir étouffant).

\end{itemize}

Généralement, lorsque les canons de protection des fosses, situés en partie haute au niveau des trémies, sont actionnés par le personnel, en salle de commande, ils permettent d’attaquer le feu de façon très précise et rapide.
\vspace{0.3cm}

Les canons à mousse, utilisés pour la protection des trémies d’enfournement, déversent très rapidement une quantité de mousse suffisante qui permet d’étouffer le feu avant que ce dernier ne crée de graves dégâts. La mousse est une quantité de goutte d’air mélangée à un liquide et son avantage, dans le cas des trémies, est que le feu est étouffé, mais surtout que le choc thermique dans le four est minime, donc il y a peu de risque de bris de l’outil, des réfractaires notamment.
\vspace{0.3cm}

Ces types d’équipements nécessitent un investissement relativement important mais permettent d’éviter de lourdes pertes en cas d’incendie.


\vspace{0.3cm}
\textbf{Rideaux d'eau ou drenchers}
\vspace{0.3cm}

Communément considérés comme synonymes, ces noms font Très souvent considérées synonymes, ces appellations désignent toutefois deux systèmes de protection différents ayant le même but. 
Le drencher (mot anglais signifiant « grosse averse ») désigne l’arrosage d’un support par de l’eau. 
\vspace{0.3cm}

Le rideau d’eau est une projection de fines gouttelettes d’eau sur une certaine longueur avec une certaine largeur formant une sorte d’écran protecteur. Les rideaux d’eau sont destinés à compléter les appareils de secours, mais en aucun cas à les remplacer ; ils ne peuvent pas non plus se substituer à un mur ou à une cloison résistant au feu.

\newpage
\vspace{0.3cm}
\textbf{Autres moyens}
\vspace{0.3cm}

Les autres moyens utilisables pour une intervention immédiate peuvent être :

\begin{itemize}
\item  des seaux d’eau ;
\item  des bacs à sable ;
\item  des couvertures anti-feu, notamment dans des locaux où l’inflammation de vêtements portés par des personnes est à craindre (atelier de mécanique).

\end{itemize}


\subsection{Matériel de deuxième intervention}


\textbf{Bouches et poteaux d’incendie}
\vspace{0.3cm}

Généralement installés à l’extérieur des locaux, les bouches et poteaux d’incendie peuvent être utilisés par le personnel, mais surtout par les pompiers qui y raccordent leur matériel. Ils doivent être incongelables, visibles et accessibles en toutes circonstances, et vérifiés périodiquement.



\subsection{Installations fixes d’extinction}


Diverses installations fixes d’extinction, généralement automatiques mais parfois manuelles, peuvent être réalisées lorsque les risques sont graves ou ponctuels (brûleurs de chaudières, stockage de produits inflammables…). Ces procédés permettent d’éteindre un foyer d’incendie par une intervention précoce et rapide, en l’absence des occupants.
